% Response_to_Reviewers.tex
\documentclass[11pt,a4paper]{article}

\usepackage[utf8]{inputenc}
\usepackage[T1]{fontenc}
\usepackage{lmodern}
\usepackage{microtype}

\usepackage{amsmath,amssymb}
\usepackage{geometry}
\usepackage{enumitem}
\usepackage{parskip}
\usepackage{hyperref}
\usepackage{xurl}
\usepackage{xcolor}

\geometry{
  top=2.5cm,
  bottom=2.5cm,
  left=2.7cm,
  right=2.7cm
}

\hypersetup{
  colorlinks=true,
  linkcolor=black,
  urlcolor=blue,
  citecolor=black
}

\setlength{\parindent}{0pt}
\setlength{\parskip}{0.7em}

% Reviewer-comment formatting
\newenvironment{reviewcomment}
  {\begin{quote}\itshape}
  {\end{quote}}

\newcommand{\response}{\noindent\textbf{Response:} }

\title{\textbf{Response to the Editor and Reviewers}}
\author{}
\date{}

\begin{document}

\maketitle
\textbf{Working response for author review; not cleared for submission.}
Submission ID: \texttt{1e58c388-d7ec-4084-b792-df20a48e3e75}.
The RMSE provenance and public data-cache requests remain open, as noted
below. Locations refer to section titles and figure/table labels rather
than page numbers from an earlier PDF.

We thank the editor and reviewers for their careful reading. We have
clarified the horizon definitions, explained their practical use, and
revised the presentation of the four-domain results. We also checked the
stored evidence behind the numerical summaries and the alternative-baseline
analysis. The corrections described below distinguish changes to the
reporting from changes to experiments. No models were fitted and no
predictions were recalculated during this editorial pass.


% ============================================================
\section*{Editor}
% ============================================================

\begin{reviewcomment}
Please report results precisely, remove overstated conclusions, and explain
limitations completely.
\end{reviewcomment}

\response
We corrected a mismatch between the PM$_{2.5}$ text and its plotted
LightGBM curve. That curve gives $H^*(\mathrm{strict})=22$ and $[27,48]$.
The former value of 13 and $[36,48]$ came from a separate moving-average
result, not the reported LightGBM curve. We corrected the text, caption,
and Table~3 without changing the forecasts or the figure data.

In ``Models'', the daily-load description now follows the producer:
expanding training, evaluation beginning after 365 observations, and a
weekly index-phase feature rather than a timestamp-derived day-of-week
label. We removed claims that the observed profiles establish intrinsic
domain limits or robustness across forecasting models. ``Limitations''
now separates positive average skill from evidence of operational benefit
and makes the baseline, metric, support, model, and horizon caps explicit.
The discussion of dynamic horizon selection remains a potential use,
not a validated controller.

The archived four-domain seasonal-baseline analysis is reported in
``Alternative-baseline sensitivity'' and Table~4. Its stored comparisons
were checked on exact common support; it was not rerun in this pass.
\textbf{Open before submission:} the inherited primary RMSE values are
not yet linked to their per-origin errors. The working manuscript does
not present them as verified results.


% ============================================================
\section*{Reviewer 1}
% ============================================================

\subsection*{Comment 1}

\begin{reviewcomment}
In the operational definitions, use a colon before the explanations of the
relaxed and strict horizons.
\end{reviewcomment}

\response
We replaced the full stop with a colon after both labels in
``Operational Predictability Horizon'' (Section~3.1).


\subsection*{Comments 2--3}

\begin{reviewcomment}
Equation (5) and the surrounding text should explicitly state that the relaxed
horizon permits non-positive skill at intermediate steps.
\end{reviewcomment}

\response
We now define
\[
\mathcal{H}^{+}
=
\{h\in\mathcal{H}:\mathrm{Skill}(h)>0\}
\]
and
\[
H^*(\mathrm{relax})
=
\max\left(\mathcal{H}^{+}\cup\{0\}\right)
\]
in Section~3.1. Equation~(5) also defines the relaxed span $\mathcal{R}$
and the gap set $\mathcal{G}=\mathcal{R}\setminus\mathcal{H}^{+}$,
whose elements have non-positive skill. This set may be non-empty.

The text explicitly notes that intermediate horizons with non-positive skill
may fall within the relaxed reach, but are not themselves interpreted as
skilful. We verified this behaviour using executable unit tests covering
intermediate skill deficits and delayed onset.


\subsection*{Comment 4}

\begin{reviewcomment}
Maintain units throughout Table 2.
\end{reviewcomment}

\response
Table~2 now specifies PM$_{2.5}$ in $\mu\mathrm{g\,m^{-3}}$, wind speed in
$\mathrm{m\,s^{-1}}$, and traffic speed in mph. For electric load, we clarify that the
stored experimental target is the sum of 15-minute kW readings and that
division by four expresses the aggregate daily demand in kWh, following the source
metadata and the aggregation script.


\subsection*{Comment 5}

\begin{reviewcomment}
Figures 1--3 are not correctly referenced in the text.
\end{reviewcomment}

\response
We have added explicit textual references to the conceptual $H^*$ schematic
(Fig.~1), the rolling-origin validation protocol (Fig.~2), and the
PM$_{2.5}$ skill curve (Fig.~3).


\subsection*{Comment 6}

\begin{reviewcomment}
Evaluate baseline choices beyond persistence.
\end{reviewcomment}

\response
We include the archived seasonal-persistence analysis, using a fixed
24-hour period for the hourly domains and a 7-day period for daily load.
For each domain, we verified equality of forecast origin, target timestamp,
horizon, and $y_{\mathrm{true}}$ between the stored LightGBM, persistence,
and seasonal-persistence records. Pairwise alignment loses no records.
These forecasts were generated in the earlier revision analysis; none
were regenerated during the present editorial pass.

For PM$_{2.5}$, the persistence comparison on this exact common support yields
$H^*(\mathrm{strict})=22$ with longest contiguous positive interval $[27,48]$,
whereas seasonal persistence yields $H^*(\mathrm{strict})=24$ with $[25,48]$.

The primary LightGBM curve also gives $H^*(\mathrm{strict})=22$ and
$[27,48]$, after the reporting correction explained above. The previous
13-hour value is not explained as a support effect. Equal descriptors
do not imply equal samples or MAEs: the seasonal-baseline comparison
uses only its matched-support rows.

Under the same matched-support analysis, load remains at
$H^*(\mathrm{strict})=1$, wind changes from 48 to 30, and traffic changes from
7 to 14. $H^*(\mathrm{relax})$ remains unchanged across all four domains. We
note that the hourly endpoints are exact multiples of the seasonal period,
where seasonal and standard persistence coincide; the unchanged relaxed
endpoint is therefore partly structural and is not interpreted as evidence of
baseline invariance.


\subsection*{Comment 7}

\begin{reviewcomment}
Explain why ARIMA(2,0,0) was selected and whether ACF was examined.
\end{reviewcomment}

\response
We agree that the original manuscript did not document this choice adequately.
An audit of our repository and Git history revealed no record that the AR(2)
order was selected via ACF/PACF, AIC/BIC, grid search, auto-ARIMA, or validation
splits. We therefore make no such claim.

The audit also revealed that the historical wind adapter requested a one-step
ARIMA forecast at every evaluated horizon, failing to implement the stated
$h$-step protocol. We removed the ARIMA numerical robustness claims from
``Models'' and the Wind and Traffic results sections rather than retaining an
unsupported rationale or modifying the historical experiment post hoc.

The separate seasonal-baseline analysis tests sensitivity to the reference
baseline, not robustness across forecasting models. We do not use it to
justify the ARIMA order or to replace evidence of model robustness. No
ARIMA fitting or order selection was performed during this revision pass.


\subsection*{Comments 8--9}

\begin{reviewcomment}
Explain the practical utility of relaxed and strict $H^*$ and what they add
beyond fixed-horizon summaries.
\end{reviewcomment}

\response
A new ``Operational interpretation and actionability'' subsection
(Section~3.2) defines $H^*(\mathrm{relax})$ as the outer horizon at
which positive baseline-relative skill is observed, and $H^*(\mathrm{strict})$,
with its interval location, as the longest contiguous positive interval.

We explain that a few fixed-horizon errors do not locate positive-skill
intervals or show the gaps between them. The PM$_{2.5}$ and traffic
examples illustrate why interval locations are needed in addition to
lengths. We also state that positive average skill alone does not establish
deployment value: its magnitude, uncertainty, decision loss, and cost matter.


\subsection*{Comment 10}

\begin{reviewcomment}
Could $H^*$ dynamically decide which horizons to forecast, and could this
create feedback or overfitting?
\end{reviewcomment}

\response
We discuss this possibility in ``Discussion''. In an
operational setting, $H^*$ could inform a decision policy estimated on
historical validation data, frozen for prospective deployment, and subsequently
updated using only realised historical errors.

We explicitly caution that estimating $H^*$, selecting horizons, and
evaluating performance on the same errors introduces adaptive selection bias.
A dynamic controller would require prospective, prequential, or nested
validation, and is not claimed to be validated here.


% ============================================================
\section*{Reviewer 2}
% ============================================================

\subsection*{Comment 1}

\begin{reviewcomment}
Explain the actionability of the findings more explicitly.
\end{reviewcomment}

\response
We added the dedicated operational interpretation subsection noted above
(Section~3.2), distinguishing the outer positive reach from the longest
contiguous positive interval and restricting dynamic deployment claims to
future, separately validated policies.


\subsection*{Comment 2}

\begin{reviewcomment}
Figures 1 and 2 are not referenced correctly.
\end{reviewcomment}

\response
Both figures are now explicitly introduced and discussed in the main text
(Figs.~1--2).


\subsection*{Comment 3}

\begin{reviewcomment}
In Figure 1, the curve overlaps text or annotations.
\end{reviewcomment}

\response
The revised Figure~1 uses dedicated callout boxes and separated label
placements. This is a layout correction, not a change to the horizon
definitions or empirical results. We inspect it in the locally compiled
manuscript; the final Overleaf rendering remains to be checked after import.


\subsection*{Comment 4}

\begin{reviewcomment}
In Figure 2, text does not fit inside the boxes.
\end{reviewcomment}

\response
The revised Figure~2 has larger boxes and wrapped text. We inspect the
locally compiled figure for overflow and clipping. The final Overleaf
rendering remains to be checked after import.


\subsection*{Comment 5}

\begin{reviewcomment}
If possible, evaluate alternative baselines beyond persistence.
\end{reviewcomment}

\response
We report the archived seasonal-persistence analysis
detailed in our response to Reviewer~1 (Comment~6), enforcing exact common
support and reporting all outcomes regardless of direction. Table~4 reports
all four domains across both baselines, documenting common-origin sample
counts, $H^*(\mathrm{relax})$, $H^*(\mathrm{strict})$, and interval locations.


\subsection*{Comment 6}

\begin{reviewcomment}
Please cache all the data in a site you control and make the central
data location clear in the paper.
\end{reviewcomment}

\response
\textbf{Not yet complete.} A README, manifest, and retrieval script help
reproduction but do not satisfy the request for a single author-controlled
public cache. The exact four-domain inputs have been identified locally.
Before submission, the permitted files must be assembled in a public
location controlled by the authors, and its downloads and checksums must
be verified. ``Data Availability'' must then point to that location and
state any redistribution restrictions. No public cache was uploaded during
this editorial pass, and we do not claim that a retrieval script is one.


\subsection*{Comment 7}

\begin{reviewcomment}
Explain the reason for ARIMA(2,0,0).
\end{reviewcomment}

\response
As noted in our response to Reviewer~1 (Comment~7), no documented
order-selection rationale was preserved, and the historical wind implementation
did not execute the intended $h$-step protocol. We removed the ARIMA robustness
claims from the text and document this constraint transparently, rather than
constructing a post hoc ACF/PACF justification.


\end{document}
